% GENERATED FILE. Edit canonical lessons, then run npm run generate:books.
% canonical-source-hash: fnv1a64:058ca11d1d4415c1
% canonical-lessons: GE-C290-bewusstlos, GE-C290-ohnmaechtig, GE-C290-bewusstsein, GE-C290-ersticken, GE-C290-ertrinken

\chapter{Bewusstlos}
\label{ch:ge-bewusstlos}
\hlchaptermodality{290}

\begin{chapteropening}
\textbf{By the end of this chapter:} \emph{I can describe a serious emergency: unconscious, faint, consciousness, to choke and to drown.}

Complete the last lesson of chapter 290: I can describe a serious emergency: unconscious, faint, consciousness, to choke and to drown.
\end{chapteropening}

\section[bewusstlos]{bewusstlos — unconscious}

\label{lesson:GE-C290-bewusstlos}



\begin{quote}
Before the new one: say the German for the fall, then the German for to bleed.
\end{quote}

\subsection*{You'll want to know: bewusstlos}
\textbf{bewusstlos} — \textquotedblleft{}unconscious\textquotedblright{}.

\begin{quote}
Der Fahrer war kurz bewusstlos.
\end{quote}

\begin{quote}
\emph{The driver was briefly unconscious.}
\end{quote}

\emph{Bewusst} 'aware' plus \emph{-los} 'without'.

\begin{cousinweb}
\emph{bewusstlos} joins \emph{bewusst} \textquotedblleft{}aware\textquotedblright{}, built on \emph{wissen} \textquotedblleft{}to know\textquotedblright{}, Proto-Germanic \emph{witaną} (EtymDB), and \emph{-los} \textquotedblleft{}without\textquotedblright{}.
\end{cousinweb}

\begin{grammarlens}[title={What you've built}]
One more word for telling, explaining and describing at B1.
\end{grammarlens}

\subsection*{Your turn}
\begin{itemize}
\raggedright
  \item \emph{Say it:} \emph{bewusstlos}
  \item \emph{Say it:} \emph{bewusstlos}, once more
  \item \emph{Say it:} a sentence of your own with \emph{bewusstlos}
  \item \emph{Recall:} say the German for the fall, then the German for to bleed, then say \emph{bewusstlos} again
\end{itemize}

\begin{tcolorbox}[breakable,colback=teal!4,colframe=teal!35!black,title={Before you move on}]
What does bewusstlos mean? (\textquotedblleft{}Unconscious\textquotedblright{}.) Say it once more.
\end{tcolorbox}
\section[ohnmächtig]{ohnmächtig — faint, passed out; powerless}

\label{lesson:GE-C290-ohnmaechtig}



\begin{quote}
Before the new one: say the German for to bleed, then the German for unconscious.
\end{quote}

\subsection*{You'll want to know: ohnmächtig}
\textbf{ohnmächtig} — \textquotedblleft{}faint, passed out; powerless\textquotedblright{}.

\begin{quote}
Sie wurde in der Hitze ohnmächtig.
\end{quote}

\begin{quote}
\emph{She fainted in the heat.}
\end{quote}

\emph{Ohnmächtig werden} — to faint. Also 'powerless'.

\begin{cousinweb}
\emph{ohnmächtig} is built on \emph{Ohnmacht}, usually traced to Middle High German \emph{āmaht} \textquotedblleft{}lack of strength\textquotedblright{}, later heard as \emph{ohne} \textquotedblleft{}without\textquotedblright{} plus \emph{Macht} \textquotedblleft{}might\textquotedblright{}, a cousin of English \emph{might} (EtymDB).
\end{cousinweb}

\begin{grammarlens}[title={What you've built}]
One more word for telling, explaining and describing at B1.
\end{grammarlens}

\subsection*{Your turn}
\begin{itemize}
\raggedright
  \item \emph{Say it:} \emph{ohnmächtig}
  \item \emph{Say it:} \emph{ohnmächtig}, once more
  \item \emph{Say it:} a sentence of your own with \emph{ohnmächtig}
  \item \emph{Recall:} say the German for to bleed, then the German for unconscious, then say \emph{ohnmächtig} again
\end{itemize}

\begin{tcolorbox}[breakable,colback=teal!4,colframe=teal!35!black,title={Before you move on}]
What does ohnmächtig mean? (\textquotedblleft{}Faint, passed out; powerless\textquotedblright{}.) Say it once more.
\end{tcolorbox}
\section[das Bewusstsein]{das Bewusstsein — consciousness; awareness}

\label{lesson:GE-C290-bewusstsein}



\begin{quote}
Before the new one: say the German for unconscious, then the German for faint.
\end{quote}

\subsection*{You'll want to know: das Bewusstsein}
\textbf{das Bewusstsein} — \textquotedblleft{}consciousness; awareness\textquotedblright{}.

\begin{quote}
Nach ein paar Minuten kam er wieder zu Bewusstsein.
\end{quote}

\begin{quote}
\emph{After a few minutes he regained consciousness.}
\end{quote}

\emph{Das Bewusstsein verlieren} — to lose consciousness.

\begin{cousinweb}
\emph{Bewusstsein} joins \emph{bewusst} \textquotedblleft{}aware\textquotedblright{}, from \emph{wissen} \textquotedblleft{}to know\textquotedblright{}, Proto-Germanic \emph{witaną} (EtymDB), and \emph{sein} \textquotedblleft{}to be\textquotedblright{}: being aware.
\end{cousinweb}

\begin{grammarlens}[title={What you've built}]
One more word for telling, explaining and describing at B1.
\end{grammarlens}

\subsection*{Your turn}
\begin{itemize}
\raggedright
  \item \emph{Say it:} \emph{das Bewusstsein}
  \item \emph{Say it:} \emph{das Bewusstsein}, once more
  \item \emph{Say it:} a sentence of your own with \emph{das Bewusstsein}
  \item \emph{Recall:} say the German for unconscious, then the German for faint, then say \emph{das Bewusstsein} again
\end{itemize}

\begin{tcolorbox}[breakable,colback=teal!4,colframe=teal!35!black,title={Before you move on}]
What does das Bewusstsein mean? (\textquotedblleft{}Consciousness; awareness\textquotedblright{}.) Say it once more.
\end{tcolorbox}
\section[ersticken]{ersticken — to choke, to suffocate}

\label{lesson:GE-C290-ersticken}



\begin{quote}
Before the new one: say the German for faint, then the German for consciousness.
\end{quote}

\subsection*{You'll want to know: ersticken}
\textbf{ersticken} — \textquotedblleft{}to choke, to suffocate\textquotedblright{}.

\begin{quote}
Der Hund erstickte fast an einem Knochen.
\end{quote}

\begin{quote}
\emph{The dog almost choked on a bone.}
\end{quote}

Perfekt with \emph{sein}: \emph{Er ist fast erstickt}. \emph{An etwas ersticken} — to choke on something.

\begin{cousinweb}
\emph{ersticken} is usually traced to Middle High German \emph{ersticken}, built on an old verb \emph{sticken} \textquotedblleft{}to get stuck, to choke\textquotedblright{}; EtymDB lists no parent.
\end{cousinweb}

\begin{grammarlens}[title={What you've built}]
One more word for telling, explaining and describing at B1.
\end{grammarlens}

\subsection*{Your turn}
\begin{itemize}
\raggedright
  \item \emph{Say it:} \emph{ersticken}
  \item \emph{Say it:} \emph{ersticken}, once more
  \item \emph{Say it:} a sentence of your own with \emph{ersticken}
  \item \emph{Recall:} say the German for faint, then the German for consciousness, then say \emph{ersticken} again
\end{itemize}

\begin{tcolorbox}[breakable,colback=teal!4,colframe=teal!35!black,title={Before you move on}]
What does ersticken mean? (\textquotedblleft{}To choke, to suffocate\textquotedblright{}.) Say it once more.
\end{tcolorbox}
\section[ertrinken]{ertrinken — to drown}

\label{lesson:GE-C290-ertrinken}



\begin{quote}
Before the new one: say the German for consciousness, then the German for to choke.
\end{quote}

\subsection*{You'll want to know: ertrinken}
\textbf{ertrinken} — \textquotedblleft{}to drown\textquotedblright{}.

\begin{quote}
Ein Mann holte das Kind aus dem Wasser, bevor es ertrank.
\end{quote}

\begin{quote}
\emph{A man pulled the child out of the water before it drowned.}
\end{quote}

Strong, with \emph{sein}: \emph{ertrank}, \emph{ist ertrunken}.

\begin{cousinweb}
\emph{ertrinken} is built on \emph{trinken} \textquotedblleft{}to drink\textquotedblright{}, from Proto-Germanic \emph{drinkaną} (EtymDB), with \emph{er-}: to drink until the end.
\end{cousinweb}

\begin{grammarlens}[title={What you've built}]
One more word for telling, explaining and describing at B1.
\end{grammarlens}

\subsection*{Your turn}
\begin{itemize}
\raggedright
  \item \emph{Say it:} \emph{ertrinken}
  \item \emph{Say it:} \emph{ertrinken}, once more
  \item \emph{Say it:} a sentence of your own with \emph{ertrinken}
  \item \emph{Recall:} say the German for consciousness, then the German for to choke, then say \emph{ertrinken} again
\end{itemize}

\begin{tcolorbox}[breakable,colback=teal!4,colframe=teal!35!black,title={Before you move on}]
What does ertrinken mean? (\textquotedblleft{}To drown\textquotedblright{}.) Say it once more.
\end{tcolorbox}
